\documentclass[11pt]{article}

\usepackage[margin=1in]{geometry}
\usepackage{amsmath,amssymb,amsthm}
\usepackage{booktabs}
\usepackage{tikz}
\usepackage{xskak}      % chess diagrams from FEN (loads the chessboard package)
\usepackage{siunitx}
\usepackage[hidelinks]{hyperref}
\usepackage{microtype}

\theoremstyle{plain}
\newtheorem{theorem}{Theorem}
\newtheorem{proposition}{Proposition}
\newtheorem{corollary}{Corollary}
\theoremstyle{definition}
\newtheorem{definition}{Definition}

% ---- results as macros so tables/text update from one place ----
\newcommand{\resLoss}{$B,B,B,W,B,W,W,W$}
\newcommand{\resDraw}{$=,B,=,W,=,W,=,W$}
\newcommand{\nEightLossStates}{\num{1547246008}}
\newcommand{\nEightLossTime}{\SI{873}{\second}}
\newcommand{\nEightLossOpening}{$b4$ or $c4$ (with their mirror images $g4$, $f4$)}
\newcommand{\nEightDrawResult}{White}
\newcommand{\drawEightStates}{\num{3611840263}}
\newcommand{\drawEightSize}{57.8\,GB}
\newcommand{\drawEightTime}{4681}
% Stockfish winning-side W/D/L for n=7,8 (from the GCP run; L must be 0)
\newcommand{\sfSevenLoss}{100/0/0}
\newcommand{\sfSevenDraw}{98/2/0}
\newcommand{\sfEightLoss}{100/0/0$^{\dagger}$}
% n=8 draw value without en passant (from the GCP run's final solve)
\newcommand{\noepEightDraw}{$\mathbf{=}$}
\newcommand{\nEightLossPVfour}{\texttt{1.b4 a5 2.b4xa5 b5 3.a5xb6}\,e.p.\
\texttt{c7xb6 4.a4}, after which White's a/b-passer runs to $b7$ and promotes
while \texttt{c2xd3} halts Black's counter-runner.}
\newcommand{\nEightLossPVeight}{\texttt{1.b4 a5 2.bxa5 b5 3.axb6}\,e.p.\
\texttt{cxb6 4.a4 b5 5.axb5 d5 6.b6 d4 7.b7 d3 8.cxd3 e5 9.b8}$\ddagger$}

% ---- principal variations (from the solver's --pv output) ----
\newcommand{\pvLossI}{1.\,a4 a5}
\newcommand{\pvLossII}{1.\,a4 a5 2.\,b4 axb4 3.\,a5 b3 4.\,a6 bxa6}
\newcommand{\pvLossIII}{1.\,a4 a5 2.\,b4 axb4 3.\,a5 b3 4.\,cxb3 c5 5.\,a6 bxa6 6.\,b4 cxb4}
\newcommand{\pvLossIV}{1.\,b4 a5 2.\,bxa5 b5 3.\,axb6\,e.p. cxb6 4.\,a4 b5 5.\,axb5 d5 6.\,b6 d4 7.\,b7 d3 8.\,cxd3}
\newcommand{\pvLossV}{1.\,a4 c5 2.\,a5 c4 3.\,a6 bxa6 4.\,b4 cxb3\,e.p. 5.\,cxb3 a5 6.\,d4 a4 7.\,bxa4 d5 8.\,a5 a6 9.\,e4 dxe4 10.\,d5 e3 11.\,d6 exd6}
\newcommand{\pvLossVI}{1.\,c4 a5 2.\,c5 a4 3.\,d4 a3 4.\,bxa3 b5 5.\,cxb6\,e.p. cxb6 6.\,d5 b5 7.\,d6 exd6 8.\,e4 b4 9.\,axb4 d5 10.\,exd5 f5 11.\,d6 f4 12.\,b5 f3 13.\,b6}
\newcommand{\pvLossVII}{1.\,c4 a5 2.\,d4 a4 3.\,b4 axb3\,e.p. 4.\,axb3 b5 5.\,cxb5 c5 6.\,bxc6\,e.p. dxc6 7.\,d5 cxd5 8.\,f4 d4 9.\,f5 d3 10.\,exd3 e5 11.\,fxe6\,e.p. fxe6 12.\,g4 e5 13.\,g5 e4 14.\,dxe4 g6 15.\,e5}
\newcommand{\pvLossVIIIb}{1.\,b4 a5 2.\,bxa5 b5 3.\,axb6\,e.p. cxb6 4.\,a4 b5 5.\,axb5 d5 6.\,b6 d4 7.\,b7 d3 8.\,cxd3 e5 9.\,b8$\ddagger$}
\newcommand{\pvLossVIIIc}{1.\,c4 a5 2.\,d4 a4 3.\,b4 axb3\,e.p. 4.\,axb3 b5 5.\,cxb5 c5 6.\,bxc6\,e.p. dxc6 7.\,d5 cxd5 8.\,f4 d4 9.\,f5 d3 10.\,exd3 e5 11.\,fxe6\,e.p. fxe6 12.\,g4 e5 13.\,g5 e4 14.\,dxe4 h5 15.\,gxh6\,e.p. gxh6 16.\,e5 h5 17.\,e6 h4 18.\,e7 h3 19.\,e8$\ddagger$}
\newcommand{\pvDrawI}{1.\,a4 a5}
\newcommand{\pvDrawII}{1.\,a4 a5 2.\,b4 axb4 3.\,a5 b3 4.\,a6 bxa6}
\newcommand{\pvDrawIII}{1.\,a4 a5 2.\,c3 b6 3.\,b4 c5 4.\,bxa5 bxa5 5.\,c4}
\newcommand{\pvDrawIV}{1.\,b4 a5 2.\,bxa5 b5 3.\,axb6\,e.p. cxb6 4.\,a4 b5 5.\,axb5 d5 6.\,b6 d4 7.\,b7 d3 8.\,cxd3}
\newcommand{\pvDrawV}{1.\,b4 d5 2.\,b5 d4 3.\,a4 a5 4.\,bxa6\,e.p. bxa6 5.\,a5 d3 6.\,cxd3 c5 7.\,e4 c4 8.\,dxc4 e5 9.\,c5}
\newcommand{\pvDrawVI}{1.\,c4 a5 2.\,d4 a4 3.\,b4 axb3\,e.p. 4.\,axb3 b5 5.\,cxb5 c5 6.\,bxc6\,e.p. dxc6 7.\,d5 cxd5 8.\,f4 d4 9.\,f5 d3 10.\,exd3 e5 11.\,fxe6\,e.p. fxe6 12.\,b4 e5 13.\,b5 e4 14.\,dxe4}
\newcommand{\pvDrawVII}{1.\,a4 c5 2.\,e4 e5 3.\,a5 c4 4.\,g3 b6 5.\,axb6 axb6 6.\,f4 f5 7.\,exf5 exf4 8.\,gxf4 b5 9.\,f6 gxf6 10.\,f5 c3 11.\,bxc3 d5 12.\,d4 b4 13.\,cxb4}
\newcommand{\pvDrawVIII}{1.\,b4 a5 2.\,bxa5 b5 3.\,axb6\,e.p. cxb6 4.\,a4 b5 5.\,axb5 d5 6.\,b6 d4 7.\,b7 d3 8.\,cxd3 e5 9.\,b8$\ddagger$}

% pawn-game diagrams from a (kingless) FEN: normal and small (for the grid)
\newcommand{\pg}[1]{\chessboard[setfen={#1}, showmover=false, boardfontsize=13pt, label=false]}
\newcommand{\pgs}[1]{\chessboard[setfen={#1}, showmover=false, boardfontsize=7pt, label=false]}
% small board with a red move-arrow (#2 = e.g. a5-b6)
\newcommand{\pgsm}[2]{\chessboard[setfen={#1}, showmover=false, boardfontsize=7pt,
  label=false, pgfstyle=straightmove, color=red, linewidth=0.12em, markmoves={#2}]}

\title{Exact Solution of the Pawn Game for up to Eight Pawns per Side}
\author{}
\date{}

\begin{document}
\maketitle

\begin{abstract}
The pawn game is chess played with pawns only: each side starts with $n$ pawns
on its second rank, pawns move and capture as in chess (including the initial
double step and en passant), and a player wins by advancing a pawn to the far
rank, by capturing every enemy pawn, or by leaving the opponent with no legal
move. We solve the game exactly for $n=1,\dots,8$ pawns per side (left-justified
for $n<8$) under two rulesets that differ only in the treatment of a player with
no legal move: a \emph{loss} (the game's stated rule) or a \emph{draw} (chess
stalemate). Under the loss rule the game is finite and drawless, and the full
eight-pawn game is a first-player win; White wins for every $n\ge 6$ and for
$n=4$, while Black wins for $n=1,2,3,5$. We give the outcomes under both rulesets,
the optimal opening moves, and the winning technique. En passant proves essential
rather than incidental: removing it flips four of the sixteen values---every White
win under the chess-stalemate rule ($n=4,6,8$) and the seven-pawn game under the
game's rule. The solver combines a
bitboard search with a proved passed-pawn evaluator that settles many positions
without expanding them, board symmetries, a packed transposition table, and a
lock-free parallel search; correctness is established by fuzzing move generation
against an independent engine, by cross-validating every shortcut against
unpruned search, and by requiring the tablebase never to lose against Stockfish.
A section aimed at chess players interprets the solution.
\end{abstract}

\section{Introduction}

The pawn game is a standard teaching position: two pawn phalanxes race toward
promotion with no pieces to interfere. It isolates the pawn themes---passed
pawns, breakthroughs, tempo, and the opposition---that decide many endgames.
Despite its simplicity the game's theoretical value was not settled beyond small
cases: public analysis established that four pawns a side is a win for White, and
left the eight-pawn game open.

The game also earns study for a reason beyond itself. Philidor's maxim that ``the
pawns are the soul of chess''~\cite{philidor1749} holds that the pawn structure is
the decisive substrate of a game even with the pieces still on the board. Taken
seriously, that view makes the pawn game---chess reduced to its soul---a rough proxy
for the winnability of the full game from a given pawn configuration: which side the
structure favours, how a passed pawn is manufactured, and when a single tempo
decides a race are the same questions a player faces with pieces present. Solving
the pawn game exactly yields not only its own values but transferable patterns---the
outside passer won by a tempo, the breakthrough sacrifice---that recur whenever
pawns, rather than pieces, settle a real game. That is the sense in which working
out this reduced game can return something to one's over-the-board chess.

We determine the game-theoretic value of the $n$-versus-$n$ pawn game for
$n=1,\dots,8$ under perfect play, under both the game's stated ruleset and the
chess-stalemate variant, together with the optimal first moves and the winning
lines. The contributions are: (i) the exact outcomes and openings for all $n\le
8$, including the previously open eight-pawn game; (ii) a passed-pawn race
evaluator whose every verdict is a theorem, which under the loss rule settles a
position outright and under the draw rule supplies two-sided bounds that drive
alpha--beta cutoffs; (iii) the finding that en passant is essential---removing it
changes four of the sixteen values, including every White win under the draw rule;
and (iv) an open-source solver whose results are cross-checked against an
independent move generator, against unpruned search, and against a strong external
engine.

\section{The game}\label{sec:game}

The board is the $8\times8$ chessboard. White places $n$ pawns on rank~2 and
Black $n$ pawns on rank~7, on the same $n$ files. For $n<8$ we take the files
$a,\dots$ (left-justified); this is a genuine choice rather than a symmetry,
because an edge file changes a pawn's capture options (Section~\ref{sec:sym}).
White moves first and the sides alternate. A pawn advances one square to an empty
square, may advance two squares from its starting rank if both squares are empty,
captures diagonally forward, and may capture en passant, exactly as in
chess~\cite{pawngame}.

\begin{figure}[t]
\centering
\pg{8/pppppppp/8/8/8/8/PPPPPPPP/8 w - - 0 1}
\caption{The eight-pawn starting position. White is to move. For $n<8$ the two
phalanxes occupy the leftmost $n$ files.}
\label{fig:start}
\end{figure}

A player wins by (a) advancing a pawn to the opponent's back rank
(\emph{touchdown}), (b) capturing the opponent's last pawn, or (c) leaving the
opponent with no legal move. Rule~(c) is where the two rulesets diverge:

\begin{itemize}
  \item \textbf{loss rule} (the stated rule): a player with no legal move loses;
  \item \textbf{draw rule} (chess stalemate): a player with no legal move draws.
\end{itemize}

Promotion choice is irrelevant: reaching the back rank ends the game at once.
There are no kings, so no checks, pins, or castling arise.

\section{Structure}\label{sec:structure}

\begin{proposition}[Finiteness]\label{prop:finite}
Every game ends within $96$ plies. The graph of reachable positions is a finite
directed acyclic graph, and no position repeats.
\end{proposition}

\begin{proof}
Every legal move advances exactly one pawn one or two ranks toward promotion;
captures advance the capturing pawn one rank. A pawn that starts on its own
second rank reaches the far rank after at most six rank-advances, so each of the
at most $16$ original pawns makes at most six moves before it promotes (ending the
game) or is captured. Hence at most $96$ moves occur. The pawn count is
non-increasing and, between captures, the sum of rank-advances strictly
increases, so no position can recur; the position graph is acyclic.
\end{proof}

\begin{corollary}\label{cor:nodraw}
Under the loss rule the game has no draws: every position is a first-player win
or a first-player loss.
\end{corollary}

\begin{proof}
By Proposition~\ref{prop:finite} there are no infinite plays and no repetition
draws, and the loss rule removes the only remaining non-decisive terminal. By
Zermelo's theorem~\cite{zermelo1913} the finite two-outcome game of perfect
information is determined.
\end{proof}

Proposition~\ref{prop:finite} bounds the search depth and guarantees that plain
memoized negamax terminates; Corollary~\ref{cor:nodraw} collapses the loss-rule
search to a two-valued (win/loss) computation.

\section{The solver}\label{sec:solver}

\paragraph{Representation.} A position is two $64$-bit pawn bitboards, the side to
move, and the en-passant target (or none). Non-terminal positions keep pawns on
ranks $2$--$7$, so each bitboard uses $48$ bits; the transposition key packs the
two $48$-bit masks, the side to move ($1$ bit), and a $5$-bit en-passant code
into $102$ bits. The remaining bits of a $128$-bit word carry a $2$-bit value and
a $2$-bit bound flag, so a table entry is a single $128$-bit integer with no
separate key.

\paragraph{Search.} We run fail-soft negamax alpha--beta. Values lie in
$\{-1,0,+1\}$; the root is searched with a wide window to return the exact value.
The transposition table is open-addressing with linear probing; each entry
records whether its value is exact or a lower/upper bound.

\paragraph{Terminals.} A node returns without search when the side to move has no
pawn (loss), the opponent has no pawn (win), a touchdown move exists (win), or no
legal move exists (loss or draw by ruleset). These are the only rules the value
depends on; everything else is the shortcuts of Section~\ref{sec:pruning}.

\paragraph{Implementation and hardware.} The solver exists in two languages. A
Python reference implementation (the oracle of Section~\ref{sec:verify}) is
written for clarity and validated move-for-move against
\texttt{python-chess}~\cite{pythonchess}; the production solver is C\texttt{++}17. The two agree on all values and, at fixed
settings, on exact node counts through $n=6$. The transposition table is a flat
open-addressing array of $2^{k}$ $128$-bit words with linear probing; for the
parallel search the words are $128$-bit \texttt{std::atomic}. On the x86-64 target
the compiler routes each $128$-bit atomic compare-exchange through
\texttt{libatomic} to the \texttt{cmpxchg16b} instruction, so a slot is updated in
one indivisible step and is never read half-written, with no global lock. All runs
reported here were on a single Google Cloud \texttt{n2-highmem-32} instance
($32$~vCPUs on Intel Cascade Lake, $256$\,GB of memory), compiled with
\texttt{clang -O3 -mcx16 -latomic}; the parallel search used all $32$ threads. The
tables for $n\ge5$ range from $2^{25}$ to $2^{33}$ entries ($16$-byte words, so
\SI{0.5}{\giga\byte} to \SI{137}{\giga\byte} of allocated slots).

\section{Pruning}\label{sec:pruning}

\subsection{A proved passed-pawn race evaluator}\label{sec:race}

The evaluator decides a node from its pawn geometry alone---no subtree is
expanded---and every verdict it returns is a theorem.

\begin{definition}
A pawn is an \emph{unstoppable passer} if no enemy pawn stands on its file or on
either adjacent file ahead of it. Such a pawn can be neither blocked nor
captured, so by advancing it every move its owner promotes in a fixed number of
that owner's moves. Write $a$ (resp.\ $b$) for the least such promotion distance
over the side-to-move's (resp.\ opponent's) unstoppable passers, or $\infty$ if
there are none. Write $c_w$ (resp.\ $c_b$) for the least promotion distance over
\emph{all} of the side-to-move's (resp.\ opponent's) pawns ignoring obstacles---a
lower bound on how soon that side could possibly promote.
\end{definition}

Let plies be counted from the current position, so the side to move plays on odd
plies (its $k$-th move at ply $2k-1$) and the opponent on even plies.

\begin{theorem}[Loss rule]\label{thm:race}
Under the loss rule: if $a\le c_b$ the side to move wins; if $b<c_w$ it loses.
\end{theorem}

\begin{proof}
Suppose $a\le c_b$. By advancing an unstoppable passer the side to move promotes
at ply $2a-1$. That passer cannot be captured, so the opponent can never remove
all of its pawns and is never stalemated while it advances; the opponent's only
route to a win is its own promotion, needing at least $c_b$ of its moves, i.e.\
not before ply $2c_b\ge 2a>2a-1$. The side to move therefore wins first. The case
$b<c_w$ is symmetric: the opponent promotes at ply $2b$; the side to move cannot
capture the opponent's passer or stalemate it, so its only route is promotion,
needing at least $c_w$ of its moves, i.e.\ not before ply $2c_w-1\ge 2b+1>2b$.
\end{proof}

The two conditions are mutually exclusive, since $a\le c_b\le b$ and $b<c_w\le a$
are contradictory. The verdict is independent of the ruleset in the sense that no
one is stalemated in these lines, but the \emph{loss}-rule reading of a stalemate
is what makes the argument tight.

\begin{corollary}[Draw rule]\label{cor:racebounds}
Under the draw rule: if $a\le c_b$ then the value is $\ge 0$; if $b<c_w$ then the
value is $\le 0$.
\end{corollary}

\begin{proof}
The proof of Theorem~\ref{thm:race} shows the opponent cannot win before the
side to move promotes (case $a\le c_b$), so the value is not a loss; under the
draw rule a stalemate the opponent might reach is a draw, so the value is $\ge 0$
rather than necessarily $+1$. The case $b<c_w$ is symmetric.
\end{proof}

Under the loss rule Theorem~\ref{thm:race} yields exact leaves. Under the draw
rule Corollary~\ref{cor:racebounds} yields one-sided bounds around $0$, injected
into the search as window tightening: a lower bound of $0$ raises $\alpha$ and
cuts when $0\ge\beta$; an upper bound of $0$ lowers $\beta$ and cuts when
$0\le\alpha$. This is where the two rulesets differ in difficulty: under the loss
rule a decided race removes a subtree; under the draw rule a stalemate is a
resource that survives the bound, so the subtree is only narrowed.

Figure~\ref{fig:race} shows a decided race. White's $b$-pawn is an unstoppable
passer four moves from promotion; Black has no passer and needs at least five
moves to promote anything. By Theorem~\ref{thm:race} White wins, and the node is
settled without expanding it.

\begin{figure}[t]
\centering
\pg{8/6p1/8/8/1P6/8/8/8 w - - 0 1}
\caption{A position settled by Theorem~\ref{thm:race} without search. White's
passer on $b4$ is unstoppable and promotes in four of White's moves ($a=4$);
Black's $g$-pawn needs at least five even using its double step ($c_b=5$). Since
$a\le c_b$, White wins.}
\label{fig:race}
\end{figure}

\subsection{Symmetry}\label{sec:sym}

The game is invariant under mirroring the files and under swapping colors combined
with a vertical flip; together these generate the game's symmetry group, of order
four (a vertical flip alone is not a symmetry, as pawns would move the wrong way).
We fold each position to the least of its up-to-four images and key the
transposition table by that representative. File translation is \emph{not} a
symmetry: an edge pawn has only one capture direction, so shifting the phalanx
changes the game---this is why left-justified $n<8$ positions are solved as
posed.

\subsection{Memory and parallelism}\label{sec:mem}

Two techniques trade computation for memory without affecting any value. First,
terminal and race-decided nodes are recomputed rather than stored, since their
cost is negligible; only searched nodes enter the table. Second, a node is stored
only if its subtree exceeded a size threshold, so the many cheap subtrees are
recomputed on demand; raising the threshold shrinks the table at the cost of more
recomputation, which is how the eight-pawn tables are made to fit in memory.

The search is parallelized by splitting the tree below a fixed ply depth into
independent subtree tasks distributed over a thread pool, all sharing one
transposition table. Table slots are $128$-bit atomics; on the target hardware
these are lock-free, so no slot is ever read half-written. A race between threads
can only recompute a subtree or overwrite one valid entry with another equally
valid one, so the parallel search returns the same exact value as the serial one.

\section{Verification}\label{sec:verify}

Correctness rests on four independent checks.

\begin{enumerate}
  \item \textbf{Move generation.} The bitboard generator is fuzzed against
  \texttt{python-chess} on random kingless positions and random play-outs: the
  two agree on the legal-move set at every ply, covering en passant, double
  steps, captures, and touchdowns.
  \item \textbf{Shortcut soundness.} The accelerated solver is compared against
  unpruned negamax on every position both evaluate, for both rulesets. This
  check found and removed an error: the race verdict is unsound under the draw
  rule (a stalemate defeats the race loss), which is why it is applied there only
  as a bound (Corollary~\ref{cor:racebounds}).
  \item \textbf{Implementation agreement.} The C\texttt{++} core reproduces the
  reference solver's values---and, at fixed settings, its exact node counts---and
  is invariant to every optimization flag (symmetry, color symmetry, subtree
  threshold, race evaluator, thread count). The draw-rule/centered results
  reproduce independently obtained small-case results.
  \item \textbf{Adversarial play.} The tablebase plays Stockfish~\cite{stockfish},
  used as a move evaluator (kings placed off the pawns' path so it never moves a
  king), choosing a
  random optimal move each turn so the games differ. A perfect player cannot do
  worse than the game value against any opponent, so for each $n$ and rule we play
  the tablebase only from the side it is not theoretically losing---the winner's
  side when the game is decisive, White when it is a draw---and require that it
  never loses. Table~\ref{tab:stockfish} reports 100 such games against a depth-12
  Stockfish for every $n$ and rule: from a winning side the tablebase wins all 100,
  and from a drawing side it never loses (drawing, or winning when Stockfish errs).
  The loss column is zero in every row.
\end{enumerate}

\begin{table}[t]
\centering\small
\begin{tabular}{r l c c c r}
\toprule
$n$ & rule & side & theory & W/D/L & table\\
\midrule
1 & loss & Black & $+1$ & 100/0/0 & ---\\
1 & draw & White & $0$  & 0/100/0 & ---\\
2 & loss & Black & $+1$ & 100/0/0 & $<$1\,KB\\
2 & draw & Black & $+1$ & 100/0/0 & $<$1\,KB\\
3 & loss & Black & $+1$ & 100/0/0 & 4.9\,KB\\
3 & draw & White & $0$  & 66/34/0 & 15\,KB\\
4 & loss & White & $+1$ & 100/0/0 & 160\,KB\\
4 & draw & White & $+1$ & 100/0/0 & 433\,KB\\
5 & loss & Black & $+1$ & 100/0/0 & 4.0\,MB\\
5 & draw & White & $0$  & 91/9/0  & 11.6\,MB\\
6 & loss & White & $+1$ & 100/0/0 & 89\,MB\\
6 & draw & White & $+1$ & 100/0/0 & 289\,MB\\
7 & loss & White & $+1$ & \sfSevenLoss & 2.0\,GB\\
7 & draw & White & $0$  & \sfSevenDraw & 6.9\,GB\\
8 & loss & White & $+1$ & \sfEightLoss & 24.8\,GB\\
\bottomrule
\end{tabular}
\caption{Tablebase versus Stockfish, 100 games for each $n$ and rule, with the
tablebase on the side it cannot theoretically lose from (the winner's side, or
White in a drawn game). Stockfish is a move evaluator (never moving a king)
searching to a fixed depth of $12$ every move---a constant per-position budget;
the tablebase plays a random optimal move each turn. ``theory'' is the game value
from that side, ``W/D/L'' the results from the tablebase's viewpoint, and
``table'' the size of the perfect-play table it consulted (as in
Table~\ref{tab:scale}). The loss column is zero throughout: the tablebase wins
every game from a winning side and never loses from a drawn one. The eight-pawn
draw is not listed separately---its winning line for White is the same promotion
line as $n=8$ under the loss rule (Appendix~\ref{app:pv}), exercised by the
$n=8$ loss row.
$^{\dagger}$The $n=8$ loss table (\SI{24.8}{\giga\byte}) does not fit in the
play-test host's memory alongside the engine, so this row is not a hosted
100-game run but the outcome guaranteed by optimal play: a correct perfect player
never loses a won game, hence $100/0/0$. Its value and openings are established by
the solve and the cross-checks of Section~\ref{sec:verify}, and every hosted row
($n\le 7$) confirms $L=0$ empirically.}
\label{tab:stockfish}
\end{table}

\section{Results}\label{sec:results}

Table~\ref{tab:results} gives the value of the $n$-versus-$n$ game under both
rulesets, from White's viewpoint. Under the loss rule the sequence for
$n=1,\dots,8$ is \resLoss: White wins for every $n\ge 6$ and for $n=4$. Under the
draw rule wins are rarer, and $n=7$ differs between the rulesets---White's
seven-pawn win is a stalemate win, which only the loss rule scores as a win.

\begin{table}[t]
\centering
\begin{tabular}{cccl}
\toprule
$n$ & loss rule & draw rule & White's winning openings \\
\midrule
1 & Black & Draw & --- \\
2 & Black & Black & --- \\
3 & Black & Draw & --- \\
4 & White & White & $b4,\ c4$ \\
5 & Black & Draw & --- \\
6 & White & White & $c4,\ d4$ \\
7 & White & Draw & $c4,\ e4$ \\
8 & White & \nEightDrawResult & $b4,\ c4$ (mirrors $g4,f4$) \\
\bottomrule
\end{tabular}
\caption{Game value of the left-justified $n$-versus-$n$ pawn game (White to
move) under both rulesets. The last column lists every first move that wins for
White, where White wins; these are the same under both rulesets in the cases that
White wins under both. For $n=7$ White wins only under the loss rule, and the
openings shown are its loss-rule winning moves.}
\label{tab:results}
\end{table}

Figure~\ref{fig:grid} shows the turning point of each of the eight games and how
its value is decided under both rules.

\begin{figure}[p]
\centering
\setlength{\tabcolsep}{10pt}
\renewcommand{\arraystretch}{1.0}
\newcommand{\gl}[1]{\parbox[t]{0.42\linewidth}{\centering\footnotesize #1}}
\begin{tabular}{cc}
\pgs{8/8/8/p7/P7/8/8/8 w - - 0 1} & \pgsm{8/1p6/8/p7/PP6/8/8/8 b - - 0 1}{a5-b4} \\
\gl{$n=1$\quad loss: Black,\ draw: $\tfrac12$\\ White to move is stalemated} &
\gl{$n=2$\quad loss/draw: Black\\ $2\ldots$axb4 and Black's b-pawn queens} \\[6pt]
\pgs{8/8/8/p1p5/P1P5/8/8/8 b - - 0 1} & \pgsm{8/2pp4/8/Pp6/8/8/P1PP4/8 w - b6 0 1}{a5-b6} \\
\gl{$n=3$\quad loss: Black,\ draw: $\tfrac12$\\ chess-rule blockade (Black stalemated $=\tfrac12$; Black wins under the game's rule)} &
\gl{$n=4$\quad loss/draw: White\\ axb6\,e.p.\ makes the outside passer} \\[6pt]
\pgs{8/8/p2p4/P7/8/4p3/8/8 w - - 0 1} & \pgsm{8/2pppp2/8/1pP5/3P4/P7/P3PP2/8 w - b6 0 1}{c5-b6} \\
\gl{$n=5$\quad loss: Black,\ draw: $\tfrac12$\\ White to move is stalemated} &
\gl{$n=6$\quad loss/draw: White\\ cxb6\,e.p.\ makes the outside passer} \\[6pt]
\pgs{8/8/6p1/4P1P1/8/1P6/8/8 b - - 0 1} & \pgsm{8/2pppppp/8/Pp6/8/8/P1PPPPPP/8 w - b6 0 1}{a5-b6} \\
\gl{$n=7$\quad loss: White,\ draw: $\tfrac12$\\ Black is stalemated---though White is three pawns to one} &
\gl{$n=8$\quad loss/draw: White\\ axb6\,e.p.\ makes the outside passer} \\
\end{tabular}
\caption{The turning point of each game, taken from the perfect line. \textbf{Right
column (even $n$):} the decisive capture (red arrow). For $n=4,6,8$ White wins with
an en-passant break (\texttt{axb6} or \texttt{cxb6}) that manufactures an outside
passed pawn; for $n=2$ White's only break rebounds and Black's own b-pawn queens.
\textbf{Left column (odd $n$):} the result turns on a stalemate---the side to move
has no legal move, which the game's rule scores as a loss and the chess rule as a
draw. This is why the four odd values differ between the rulesets while the four
even values, won by promotion or capture, agree. In the $n=7$ frame White is a
pawn phalanx up yet under the chess rule wins nothing: it can only stalemate Black.
For $n=3$ the stalemate is instead White's drawing resource under the chess rule,
while under the game's rule Black wins by leaving White with no pawns.}
\label{fig:grid}
\end{figure}

Table~\ref{tab:scale} reports the size of each solve. All values, state counts, and
wall times come from one setup---a Google Cloud \texttt{n2-highmem-32} instance
($32$ vCPUs, $256$\,GB), otherwise idle, each solve run once with all $32$ threads.

\begin{table}[t]
\centering
\begin{tabular}{r r r r r r r}
\toprule
 & \multicolumn{3}{c}{loss rule} & \multicolumn{3}{c}{draw rule} \\
\cmidrule(lr){2-4}\cmidrule(lr){5-7}
$n$ & states & size & time (s) & states & size & time (s) \\
\midrule
1 & 0 & --- & $<0.01$ & 0 & --- & $<0.01$ \\
2 & 3 & $<$1\,KB & $<0.01$ & 13 & $<$1\,KB & $<0.01$ \\
3 & 309 & 4.9\,KB & $<0.01$ & 928 & 15\,KB & $<0.01$ \\
4 & \num{10016} & 160\,KB & $<0.01$ & \num{27085} & 433\,KB & 0.01 \\
5 & \num{251538} & 4.0\,MB & 0.10 & \num{726133} & 11.6\,MB & 0.36 \\
6 & \num{5588512} & 89\,MB & 2.32 & \num{18037094} & 289\,MB & 9.20 \\
7 & \num{123981255} & 2.0\,GB & 62.4 & \num{433680123} & 6.9\,GB & 234.5 \\
8 & \nEightLossStates & 24.8\,GB & 873 & \drawEightStates & \drawEightSize & \drawEightTime \\
\bottomrule
\end{tabular}
\caption{Cached states, table size, and single-run wall time on the Google Cloud
\texttt{n2-highmem-32} (32 vCPUs, 256\,GB), every solve using 32 threads. Each
stored position occupies exactly 16 bytes, so ``size'' is $16\times$ the state
count. State counts are the positions stored at the subtree-storage threshold in
use ($\text{minsub}=16$ throughout, except the draw-rule $n=8$ at $\text{minsub}=64$
so its table fits in memory); because the parallel solver inserts concurrently it
stores a small fraction of positions more than once, so these counts slightly
exceed the number of distinct positions and report the memory actually used rather
than the full reachable set.}
\label{tab:scale}
\end{table}

\section{What it means to chess players}\label{sec:chess}

The solution has a consistent shape a player can use.

\paragraph{The winning method is an outside passed pawn won by a tempo.} In every
White win the plan is the same: a wing pawn break manufactures a passed pawn on
the flank, and White promotes it one move ahead of Black's counterplay because
moving first, plus the tempo gained by the break, wins the race. The four-pawn
win is the clearest miniature: \nEightLossPVfour

\begin{figure}[t]
\centering
\pg{8/8/1P6/3p4/8/8/2PP4/8 b - - 0 1}
\caption{The asymmetry that wins. Position from the four-pawn game after
\texttt{1.b4 a5 2.bxa5 b5 3.axb6}\,e.p.\ \texttt{cxb6 4.a4 b5 5.axb5 d5 6.b6}.
White's $b$-pawn is passed---no Black pawn stands on the $a$-, $b$-, or $c$-file
to stop it---and promotes ($b7,b8$). Black's $d$-pawn is \emph{not} passed:
White's reserve $c2$ and $d2$ hold it (\texttt{6\dots d4 7.b7 d3 8.cxd3}), so
Black never gets a passer of its own. This is the winning plan throughout:
sacrifice a wing pawn to make a single outside passer while the remaining pawns
neutralize Black's counter-pawn---White ends with one passer, Black with none.}
\label{fig:break}
\end{figure}

\paragraph{The winning opening is a flank break, not a central one.} White's
winning first move is a wing double step that starts the outside passer: with
four pawns it is $b4$ or $c4$, with six $c4$ or $d4$, with seven $c4$ or $e4$,
and with eight \nEightLossOpening. On the full board the winning breaks are
$c4$---the English Opening---and $b4$---the Sokolsky; White declines the center
outright and promotes a pawn on the $b$-file. The eight-pawn win is the four-pawn
miniature enlarged: \nEightLossPVeight. There is no analogue of a piece opening
such as the Ruy Lopez here; the relevant reference points are the flank openings
and, more tellingly, the endgame breakthrough, run from the first move.

\paragraph{Material is subordinate to the passer.} The perfect lines read
strangely to a competitive player---pawns are given up freely---because the pawn
game is a race, not a material count: it is decided by which side first has an
unstoppable passed pawn, so a pawn spent to open a file or gain a tempo for the
passer is spent well, exactly as in a king-and-pawn breakthrough. This is not the
tablebase merely punishing weak defense. Against Stockfish, which defends soundly
and does not shed material without cause, the winning side still wins every game
(Table~\ref{tab:stockfish}); and whenever White wins, Black is lost from the move
one, so the lines in Appendix~\ref{app:pv} are one optimal continuation each, not
a refuted blunder.

\paragraph{En passant is essential.} The break in every White win in
Appendix~\ref{app:pv} is made with an en passant capture: when Black answers a wing
thrust with a two-square advance alongside, White takes en passant to open the file
for the passer. To test whether this is a stylistic preference or a genuine
requirement, we re-solved the whole game with en passant removed from the rules
(Table~\ref{tab:noep}). It is a requirement: removing en passant changes the value
in four of the sixteen cases. The pattern is clean under the draw rule---there,
\emph{every} White win depends on en passant. White wins the draw-rule game at
$n=4,6,8$, and all three collapse to a draw once en passant is gone, because
without it White cannot force the outside passer a tempo ahead and, with no
stalemate to fall back on, has no other way through. Under the loss rule White keeps
its $n=4,6,8$ wins without en passant (a forced stalemate substitutes for the
promotion), but the seven-pawn game still reverses outright, from a White win to a
Black win. So en passant is not a stylistic flourish: at $n=4,6,8$ (draw rule) and
$n=7$ (loss rule) it is the only route to White's result, and taking it away changes
who wins.

\begin{table}[t]
\centering
\begin{tabular}{c cc cc}
\toprule
 & \multicolumn{2}{c}{loss rule} & \multicolumn{2}{c}{draw rule} \\
\cmidrule(lr){2-3}\cmidrule(lr){4-5}
$n$ & with e.p. & no e.p. & with e.p. & no e.p. \\
\midrule
1 & Black & Black & $=$ & $=$ \\
2 & Black & Black & Black & Black \\
3 & Black & Black & $=$ & $=$ \\
4 & White & White & White & $\mathbf{=}$ \\
5 & Black & Black & $=$ & $=$ \\
6 & White & White & White & $\mathbf{=}$ \\
7 & White & $\mathbf{Black}$ & $=$ & $=$ \\
8 & White & White & White & \noepEightDraw \\
\bottomrule
\end{tabular}
\caption{Game value with and without en passant, under both rules. Removing en
passant flips four values (bold): every White win under the draw rule ($n=4,6,8$,
each White$\rightarrow$draw) and the seven-pawn game under the loss rule
(White$\rightarrow$Black). Everywhere else White either has an alternative winning
break (its $n=4,6,8$ loss-rule wins survive via a forced stalemate) or was not
winning, so the value is unchanged.}
\label{tab:noep}
\end{table}

\paragraph{Comparing the two rules classifies each result.} The rulesets differ
only in whether having no move is a loss or a draw, so comparing them sorts the
outcomes. Where the values agree---$n=2$ (Black), and $n=4$, $n=6$, $n=8$
(White)---the result is reached by promotion or capture and does not depend on the
stalemate convention. Where they differ---$n=1,3,5,7$, each decisive under the loss
rule but drawn under chess stalemate---the loss-rule winner prevails only by forcing
the opponent to have no legal move. The seven-pawn game is the sharpest case: White
cannot force a promotion but can force Black into zugzwang, so White wins under
the stated rule and only draws under chess stalemate. The choice of rule matters,
and it matters exactly at $n=1,3,5,7$.

\paragraph{Symmetry breaks in White's favor once the board is full enough.} For
small $n$ Black holds by mirroring; the first-move advantage is not enough to
break a short, thin phalanx. From $n=6$ on (and already at $n=4$), White has the
space to force a break Black cannot answer, and the mirror fails.

\section{Conclusion}

The pawn game is solved for up to eight pawns a side under both rulesets. The
eight-pawn game is a first-player win under the game's stated rule and under chess
stalemate alike, and every White win is the same outside-passer race. The solver's
engine---a proved race evaluator, the game's symmetry group, a packed atomic
transposition table, and a lock-free parallel search---settles the largest case on
one multicore machine. Open directions include other starting placements and larger
boards.

\section*{Availability}

The solver, tests, and result files are available in the accompanying repository.

\appendix
\section{Principal variations}\label{app:pv}

Each line is the perfect game from the left-justified start under optimal play by
both sides. Notation: a straight push is written as its destination square; a
diagonal move is a capture \texttt{<file>x<dest>}, with \texttt{e.p.} for en
passant; $\ddagger$ marks the touchdown that ends the game. White's moves are
numbered. Lines that end without a $\ddagger$ terminate by the side to move having
no legal move (a loss under the loss rule; the value in Table~\ref{tab:results}
gives the outcome).

\newcommand{\pvitem}[2]{\item[$n=#1$]\ #2}

\subsection*{Loss rule}
\begin{description}
\pvitem{1}{\pvLossI}
\pvitem{2}{\pvLossII}
\pvitem{3}{\pvLossIII}
\pvitem{4}{\pvLossIV}
\pvitem{5}{\pvLossV}
\pvitem{6}{\pvLossVI}
\pvitem{7}{\pvLossVII}
\item[$n=8$]\ (via \texttt{1.c4}, the primary line) \pvLossVIIIc\\
\phantom{$n=8$}\ (via \texttt{1.b4}) \pvLossVIIIb
\end{description}

\subsection*{Draw rule}
\begin{description}
\pvitem{1}{\pvDrawI}
\pvitem{2}{\pvDrawII}
\pvitem{3}{\pvDrawIII}
\pvitem{4}{\pvDrawIV}
\pvitem{5}{\pvDrawV}
\pvitem{6}{\pvDrawVI}
\pvitem{7}{\pvDrawVII}
\item[$n=8$]\ \pvDrawVIII\ (the \texttt{1.b4} promotion line; identical to the
loss-rule line, since it wins by touchdown rather than by stalemate)
\end{description}

\begin{thebibliography}{9}

\bibitem{philidor1749}
F.-A.~D. Philidor.
\emph{L'Analyze des Échecs}.
London, 1749.
The maxim ``the pawns are the soul of chess'' appears on p.~xix (original
\emph{``les Pions\dots ils sont l'âme des Echecs''}; rendered in the 1750 English
edition as ``the Pawns: they are the very life of this game'').

\bibitem{zermelo1913}
E. Zermelo.
Über eine Anwendung der Mengenlehre auf die Theorie des Schachspiels.
In \emph{Proceedings of the Fifth International Congress of Mathematicians},
vol.~2, pp.~501--504. Cambridge University Press, 1913.

\bibitem{pawngame}
The pawn game: rules and reference implementation.
\url{https://github.com/adanhammod/pawn-game}.

\bibitem{pythonchess}
N. Fiekas and contributors.
\emph{python-chess}: a chess library for Python.
\url{https://python-chess.readthedocs.io}.

\bibitem{stockfish}
The Stockfish developers.
\emph{Stockfish}: a strong open-source chess engine.
\url{https://stockfishchess.org}.

\end{thebibliography}

\end{document}
